% SOURCE_AUTHORITY: sga5_fr_workpass.tex, lines 13549--13576 (LF-normalized unit).
% SOURCE_SHA256: 791F4EFFC5E02832D5D77ED03518C8156D6F07E4C8238B03545DB93D883FBB28
O homomorfismo $u(f'_R)$ é, de fato, um homomorfismo de espaços vetoriais sobre $k$. Para vê-lo, é preciso provar que, quaisquer que sejam $x \in \fm$, $\alpha \in R$,
\[
f'_R(\alpha x) - \alpha x \equiv \alpha(f'_R(x) - x) \bmod \fm^{r+1} \quad ,
\]
Isto é,
\begin{equation}
\tag{4.7.1}
f'_R(\alpha x)-\alpha f'_R(x)\in\fm^{r+1}.
\end{equation}
Mas
\[
f'_R(\alpha x)-\alpha f'_R(x)=(f'_R(\alpha)-\alpha)f'_R(x),
\]
o que fornece a relação (4.7.1), porque
\[
f'_R(\alpha)-\alpha\in\fm^r\quad\hbox{(condição a))},
\qquad
f'_R(x)\in\fm\quad (x\in\fm).
\]
O homomorfismo $u(f'_R)$ é um homomorfismo de $k[G]$-módulos. De fato, para $x\in\fm$, tem-se:
\[
\begin{aligned}
u(f'_R)(g(x\bmod\fm^2))
&=(f'_R(gx)-gx)\bmod\fm^{r+1}\\
&=g(f'_R(x)-x)\bmod\fm^{r+1}
=gu(f'_R)(x\bmod\fm^2).
\end{aligned}
\]
